\section*{Chapitre \ref{ch:iv1:timed-test-10} --- Épreuve chronométrée 10}
\markboth{\protect\omTocMark{}{Réponses~: Épreuve chronométrée 10}}{\protect\omTocMark{}{Réponses~: Épreuve chronométrée 10}}

\begin{solution}{iq:iv1:timed-test-10:1}
La probabilité vaut $\frac23$.

La somme est inférieure à 1 dans le simplexe de coin du cube, dont le volume vaut $\frac{1}{3!} = \frac16$\index{volume d'un simplexe}. Par la symétrie $x \to 1 - x$ appliquée aux trois variables, la somme dépasse 2 avec la même probabilité. Le milieu a donc pour probabilité $1 - \frac26 = \frac23$.

Une simulation donne $0.667$. Ici, la moyenne vaut $1.5$ et l'écart-type $\frac12$, de sorte que l'intervalle s'étend d'un écart-type de part et d'autre, et $\frac23$ est proche de la règle empirique de la loi normale, soit $0.68$.

\iqwrong{Traiter la somme comme uniforme sur $[0, 3]$, ce qui donne $\frac13$. Les sommes se concentrent vers le centre.}
\end{solution}

\begin{solution}{iq:iv1:timed-test-10:2}
La probabilité vaut $\frac{50}{81} \approx 0.617$.

Il y a $3^5 = 243$ issues. Utiliser l'inclusion--exclusion\index{inclusion-exclusion} sur les cases vides. Pour une case donnée vide, il y a $2^5 = 32$ issues, et pour deux cases données vides, il n'y en a qu'une. Les issues sans case vide sont donc
\[
243 - 3 \times 32 + 3 \times 1 = 150.
\]
Cela donne $\frac{150}{243} = \frac{50}{81}$.

Un second dénombrement utilise les tailles~: $(3, 1, 1)$ compte $3 \times \binom53 \times 2 = 60$ issues, et $(2, 2, 1)$ en compte $3 \times \binom52\binom32 = 90$. Le total vaut 150. Une énumération donne le même résultat.

\iqwrong{Oublier le dernier terme et obtenir $\frac{147}{243}$. Les issues à deux cases vides sont retranchées deux fois dans le comptage à une case vide~: il faut donc les rajouter.}
\end{solution}

\begin{solution}{iq:iv1:timed-test-10:3}
L'inégalité de Tchebychev donne au plus $0.16$, et la valeur pour la loi normale est d'environ $0.0124$.

L'inégalité de Tchebychev affirme que $P(|X - \mu| \ge a) \le \frac{\sigma^2}{a^2}$\index{inégalité de Tchebychev}, soit $\frac{4}{25} = 0.16$.

Pour une loi normale, l'écart vaut $\frac52 = 2.5$ écarts-types. Les deux queues comptent~: $2 \times 0.0062 = 0.0124$. Une simulation confirme ce résultat.

L'écart, d'un facteur 13, est le prix de l'absence de toute hypothèse sur la forme. Et la borne de Tchebychev est atteinte par une loi à trois points, qui place toute la déviation exactement en $\pm5$. Elle ne peut donc pas être améliorée sans hypothèse supplémentaire.

Les deux valeurs sont justes~: seule la première vaut pour toute loi de cette variance, et cela compte lorsqu'une queue est épaisse.

\iqwrong{Ne prendre qu'une seule queue pour la loi normale, soit $0.0062$. La question porte sur un écart dans un sens ou dans l'autre.}
\end{solution}

\begin{solution}{iq:iv1:timed-test-10:4}
Chaque probabilité vaut $\frac59$.

Chaque dé a trois valeurs, chacune avec la probabilité $\frac13$. Il y a donc neuf couples équiprobables.

A contre B~: 2 bat 1, 4 bat 1, et 9 bat les trois. Cela fait $1 + 1 + 3 = 5$ couples sur 9.

B contre C~: 6 bat 3 et 5, et 8 bat les trois. Cela fait $2 + 3 = 5$.

C contre A~: 3 bat 2, 5 bat 2 et 4, et 7 bat 2 et 4. Cela fait $1 + 2 + 2 = 5$.

Ainsi A bat B, B bat C et C bat A. Une énumération donne le même résultat. La relation «~bat~» n'est pas transitive\index{relation non transitive}, bien que chaque dé ait la même moyenne, soit 5. Celui qui choisit en second a donc toujours l'edge.

\iqwrong{Supposer qu'un dé doit être le meilleur, parce qu'un classement paraît naturel. Classer par probabilité de victoire en duel ne donne pas forcément un ordre.}
\end{solution}

\begin{solution}{iq:iv1:timed-test-10:5}
La probabilité vaut $\frac23$, et l'attente espérée vaut $\frac{20}{3} \approx 6.67$ minutes.

Les taux valent $\frac{1}{10}$ et $\frac{1}{20}$. Par la propriété d'absence de mémoire, l'attente pour chacun suit une loi exponentielle, quoi qu'il soit arrivé avant votre arrivée. Le premier bus est sur la ligne A avec la probabilité\index{course exponentielle} $\frac{1/10}{1/10 + 1/20} = \frac23$.

Le premier passage de bus a pour taux $\frac{3}{20}$, donc son espérance vaut $\frac{20}{3}$. Une simulation confirme ce résultat.

Vérification aux limites~: si la ligne B est supprimée, son taux vaut 0, la probabilité vaut 1 et l'attente vaut 10. L'attente est donc inférieure à l'écart le plus court, ce qui est cohérent.

Le choix de la ligne et le temps d'attente sont indépendants. L'attente vaut donc $\frac{20}{3}$ quelle que soit la ligne que l'on finit par prendre.

\iqwrong{Prendre la moitié de chaque écart comme attente, ce qui n'est juste que pour des bus réguliers. Sur une ligne suivant un processus de Poisson, l'attente moyenne est égale à l'écart moyen complet.}
\end{solution}

\begin{solution}{iq:iv1:timed-test-10:6}
Le plus petit groupe compte 6 personnes.

La probabilité que les $n$ anniversaires soient tous différents\index{problème des anniversaires} est $\prod_{i = 0}^{n - 1}(1 - \frac{i}{20})$. Pour $n = 5$, elle vaut $\frac{20 \cdot 19 \cdot 18 \cdot 17 \cdot 16}{20^5} = 0.581$. Pour $n = 6$, elle vaut $0.581 \times \frac{15}{20} = 0.436$. Ainsi la probabilité d'une coïncidence passe de $0.419$ à $0.564$, et elle dépasse $\frac12$ dès 6.

La règle empirique est $1.18\sqrt{d}$~: avec $d = 20$, elle donne $5.3$, qu'on arrondit à 6. Un calcul exact par boucle le confirme. La même règle donne $22.5$ pour 365 jours, proche de la valeur exacte, 23.

\iqwrong{Répondre 10, soit la moitié des jours, parce que la moitié paraît naturelle. Les paires croissent comme $n^2$, si bien que la réponse évolue comme la racine carrée du nombre de jours.}
\end{solution}

\begin{solution}{iq:iv1:timed-test-10:7}
L'espérance vaut $\frac{0.5}{\ln 2} \approx 0.721$.

La densité conjointe vaut $\frac1x$ pour $0 < y < x < 1$. Ainsi, sachant $Y = y$, la densité de $X$ est proportionnelle à $\frac1x$ sur $(y, 1)$\index{densité conditionnelle}. La constante vaut $-\ln y$. Alors
\[
\mathbb{E}[X \mid Y = y] = \frac{\int_y^11\,dx}{-\ln y} = \frac{1 - y}{-\ln y}.
\]
Pour $y = \frac12$, cela donne $\frac{0.5}{0.693} = 0.721$.

Vérification~: elle est comprise entre $\frac12$ et 1, comme il se doit. Lorsque $y \to 1$, la formule tend vers 1. Lorsque $y \to 0$, elle tend vers 0, ce qui est juste, car un petit $Y$ indique un petit $X$. Une simulation dans une fine fenêtre autour de $\frac12$ le confirme.

\iqwrong{Répondre $\frac34$, le milieu de l'intervalle admissible. Le facteur $\frac1x$ favorise les petits $X$, car un petit $X$ porte davantage de densité pour un $Y$ donné.}
\end{solution}

\begin{solution}{iq:iv1:timed-test-10:8}
Avec cinq dés, la probabilité vaut exactement $\frac12$, et avec six dés elle vaut environ $0.546$.

Remplacer chaque face $x$ par $7 - x$, ce qui laisse inchangée la loi de chaque dé\index{symétrie}. Avec cinq dés, le total $T$ devient $35 - T$, de sorte que $T$ et $35 - T$ ont la même loi. Or $T \ge 18$ est le même événement que $35 - T \le 17$. Donc $P(T \ge 18) = P(T \le 17)$, et ces deux événements recouvrent tout, sans intersection. Chacun a donc pour probabilité $\frac12$.

Avec six dés, $T$ et $42 - T$ ont la même loi, et le centre, 21, est un total possible. Donc $P(T > 21) = P(T < 21)$, et
\[
P(T \ge 21) = \frac12 + \frac12P(T = 21) = \frac12 + \frac12 \cdot \frac{4\,332}{46\,656} \approx 0.546.
\]

\iqwrong{Donner $\frac12$ aussi pour six dés. Cinq dés ont un centre égal à $17.5$, valeur qu'aucun total ne peut prendre, mais six dés ont un centre atteignable, et «~au moins~» l'inclut.}
\end{solution}

\begin{solution}{iq:iv1:timed-test-10:9}
Le troisième élément est conservé avec la probabilité $\frac15$. Pour un échantillon de taille 2, on conserve chaque nouvel élément avec la probabilité $\frac2i$, en remplaçant au hasard l'un des deux, et chaque élément a alors la probabilité $\frac25$.

Le troisième élément entre avec la probabilité $\frac13$. Il survit au quatrième avec $\frac34$ et au cinquième avec $\frac45$\index{échantillonnage par réservoir}. Le produit vaut $\frac13 \cdot \frac34 \cdot \frac45 = \frac15$.

Le produit se télescope pour tous $i$ et $n$, de sorte que chaque élément finit avec la probabilité $\frac1n$. On n'a jamais besoin de connaître $n$ à l'avance.

Pour un échantillon de taille $k$, l'élément $i > k$ entre avec la probabilité $\frac{k}{i}$. Chaque élément ultérieur $j$ l'élimine avec la probabilité $\frac{k}{j} \cdot \frac1k = \frac1j$, de sorte qu'il survit avec $\frac{j - 1}{j}$. Le produit vaut $\frac{k}{n}$. Une simulation confirme que les parts sont uniformes.

\iqwrong{Utiliser une probabilité fixe par élément, disons $\frac{1}{100}$. Cela suppose de connaître la longueur, et cela avantage les éléments ultérieurs.}
\end{solution}

\enlargethispage{-6\baselineskip}%
\begin{solution}{iq:iv1:timed-test-10:10}
La probabilité vaut environ $0.0211$.

Choisir le rang~: 13. Choisir ses trois couleurs\index{dénombrement des mains}~: $\binom43 = 4$. Choisir les deux autres rangs~: $\binom{12}{2} = 66$. Choisir une couleur pour chacun~: $4^2 = 16$. Le compte vaut
\[
13 \times 4 \times 66 \times 16 = 54\,912,
\]
sur le total, soit $0.0211$.

Une simulation indépendante renvoie la même valeur. Notons que les deux rangs doivent être différents, sinon la main serait un full.

\iqwrong{Utiliser 49 choix pour la quatrième carte et 48 pour la cinquième. Cela compte les full, et chaque main est comptée deux fois.}

\iqfollow{Quelle est la probabilité d'un full~?}{$13 \times 4 \times 12 \times 6 = 3744$ mains, soit $0.00144$. Il est plus rare, bien que la description soit semblable.}
\end{solution}

\begin{solution}{iq:iv1:timed-test-10:11}
La probabilité vaut environ $0.057$.

Le compte a pour espérance 50 et pour écart type 5\index{correction de continuité}. On veut $H \le 40$ ou $H \ge 60$. La correction place les bornes en $40.5$ et $59.5$, chacune à $1.9$ écart type de la moyenne. Les deux queues donnent $2 \times 0.0287 = 0.0574$.

Une somme binomiale exacte donne $0.0569$, la correction est donc très proche. Sans correction, $z = 2$, et la réponse serait $0.0455$. L'écart est d'un cinquième, une erreur importante pour une étape aussi simple.

Une vérification avec un chiffre connu~: deux écarts types couvrent environ $95~\%$, donc environ $5~\%$ se situent au-delà. Un compte qui peut tomber exactement sur la borne ajoute un peu plus.

\iqwrong{Utiliser une seule queue, ce qui donne $0.029$, soit la moitié de la réponse. La question demande un écart dans les deux sens.}
\end{solution}

\begin{solution}{iq:iv1:timed-test-10:12}
Le put devrait coûter 15. À 14, acheter le put et l'action, et vendre le call, pour un profit certain de 1.

La parité call-put dit que $C - P = S - K$ à taux nul\index{parité call-put}. Donc $P = C - S + K = 10 - 100 + 105 = 15$.

À 14, le put est trop bon marché. Acheter le put, acheter l'action et vendre le call. Le coût vaut $14 + 100 - 10 = 104$. À l'échéance, l'action plus le put moins le call valent exactement $K = 105$ dans tous les cas~: si $S_T > 105$, le call est exercé contre vous à 105, et si $S_T < 105$, vous exercez le put à 105. Le profit est donc de 1, sans risque.

Une vérification~: pour $S_T = 90$ et pour $S_T = 120$, la position vaut 105.

\iqwrong{Vendre le put, parce qu'il paraît mauvais à 14. La règle consiste à acheter ce qui est bon marché par rapport à sa réplication.}
\end{solution}
